pdfoutput=1
\pdfoutput=1
\documentclass[12pt,a4paper]{article}

% ============================================================
%  FA-Theorem — Federated Audit Theorem:
%  Zero-Knowledge Multi-Party Audit Completeness via
%  Situs Homomorphic Composition
%  arXiv-ready · English · Standalone (SCX-citable)
% ============================================================

\usepackage[T1]{fontenc}
\usepackage{amsmath,amssymb,amsfonts}
\usepackage{mathtools}
\usepackage{amsthm}
\usepackage{bm}
\usepackage{graphicx}
\usepackage{xcolor}
\usepackage{hyperref}
\usepackage{booktabs}
\usepackage{enumitem}
\usepackage[numbers]{natbib}
\usepackage{algorithm}
\usepackage{algpseudocode}

% ---------- Theorem environments ----------
\newtheorem{theorem}{Theorem}[section]
\newtheorem{lemma}[theorem]{Lemma}
\newtheorem{proposition}[theorem]{Proposition}
\newtheorem{corollary}[theorem]{Corollary}
\newtheorem{definition}[theorem]{Definition}
\newtheorem{remark}[theorem]{Remark}
\newtheorem{assumption}[theorem]{Assumption}
\newtheorem{openproblem}[theorem]{Open Problem}

% ---------- Status tags ----------
\newcommand{\rigorous}{\textcolor{green!50!black}{\small\textsf{[Rigorous]}}}
\newcommand{\conditionallyrigorous}{\textcolor{orange}{\small\textsf{[Conditionally Rigorous]}}}
\newcommand{\openquest}{\textcolor{red!70!black}{\small\textsf{[Open]}}}

% ---------- Macros ----------
\newcommand{\R}{\mathbb{R}}
\newcommand{\E}{\mathbb{E}}
\newcommand{\V}{\mathbb{V}}
\newcommand{\I}{\mathbb{I}}
\newcommand{\cX}{\mathcal{X}}
\newcommand{\cY}{\mathcal{Y}}
\newcommand{\cD}{\mathcal{D}}
\newcommand{\cS}{\mathcal{S}}
\newcommand{\cP}{\mathcal{P}}
\newcommand{\cA}{\mathcal{A}}
\newcommand{\cI}{\mathcal{I}}
\newcommand{\ind}[1]{\mathbf{1}_{[#1]}}
\newcommand{\argmax}{\operatorname*{arg\,max}}
\newcommand{\argmin}{\operatorname*{arg\,min}}
\newcommand{\KL}{D_{\mathrm{KL}}}
\newcommand{\Situs}{\mathrm{Situs}}
\newcommand{\Cercis}{\mathrm{Cercis}}
\newcommand{\Commit}{\mathrm{Commit}}
\newcommand{\Verify}{\mathrm{Verify}}
\newcommand{\ZK}{\mathrm{ZK}}
\newcommand{\negl}{\mathrm{negl}}
\newcommand{\poly}{\mathrm{poly}}
\newcommand{\Power}{\mathrm{Power}}
\newcommand{\dSitus}{d_{\mathrm{Situs}}}
\DeclareMathOperator{\Tr}{Tr}
\DeclareMathOperator{\diag}{diag}
\newcommand{\Wass}{\mathcal{W}}

\title{\bfseries Federated Audit Completeness:\\
Zero-Knowledge Multi-Party Auditing\\
via Situs Homomorphic Composition}

\author{SCX}

\date{\today}

\begin{document}
\maketitle

\begin{abstract}
The \textsf{SCX} Audit Sword establishes that any use of the Cercis Score can be
independently audited --- provided the auditor possesses complete data access.
In federated scenarios where $K$ parties hold private datasets and interact
only through aggregated statistics, the Audit Sword encounters a structural
obstacle: how does one audit what one cannot see?  We formalize this as the
\emph{Federated Audit Problem} and prove three results.  (i)~When the Situs
encoding admits a homomorphic composition operator $\oplus$ such that
$\Situs(\cup_k \cD_k) \cong \oplus_k \Situs(\cD_k)$, and each party's Cercis
Score submission is verified by a zero-knowledge proof $\pi_{\ZK}$, the
federated audit achieves $\varepsilon$-equivalence to the full-data audit,
with $\varepsilon \to 0$ under vanishing differential privacy budget and
growing sample size.  The Situs homomorphic composition construction is
\textbf{conditionally rigorous} --- given the operator, the $\varepsilon$-bound
is strict; the operator's existence is an \textbf{open problem}.
(ii)~Without $\ZK$ verification, the information loss is lower-bounded by
$\mathrm{const} \cdot \eta_k \cdot |\cD_k|/|\cD_{\mathrm{total}}|$, a
\textbf{rigorous} consequence of Theorem~3 and the data processing inequality.
(iii)~We outline the ``Telegraph Public Good'' extension where decentralized
storage-layer commitments replace $\ZK$ proofs with Proofs of Storage,
further tightening the $\varepsilon$ bound --- an \textbf{open problem} with
a clear theoretical path.  Collectively, the results establish that federated
auditing is possible without raw data access, but only when the information
architecture provides verifiable audit fingerprints.
\end{abstract}

% ============================================================
\section{Introduction}
% ============================================================

The \textsf{SCX} framework~\citep{scx2024cercis} established the \textbf{Audit Sword}
principle: every use of the Cercis Score must be independently verifiable by
any third party.  This is not an optional feature --- it is a structural
requirement that follows from Theorem~3 (the Noise--Difficulty
Unidentifiability Theorem).  If you cannot audit the detector, you cannot
distinguish whether it is flagging genuine label noise or merely reclassifying
difficult-but-correct samples according to its own inductive biases.

The Audit Sword, however, carries a hidden premise: that the auditor has
\textbf{complete access to all data}.  In an increasingly federated world ---
hospitals that cannot share patient records across jurisdictions, financial
institutions barred from exposing transaction logs, edge devices with
bandwidth and privacy constraints --- this premise fails.  The very parties
whose models most urgently require auditing (because they operate on
sensitive, high-stakes data) are precisely those who cannot grant the access
that auditing presupposes.

This creates the \textbf{Federated Audit Paradox}.  Federated learning has
developed an extensive privacy-preserving toolbox --- Secure
Aggregation~\citep{bonawitz2017practical}, Differential
Privacy~\citep{abadi2016deep}, Multi-Party
Computation~\citep{evans2018pragmatic} --- but each of these tools is designed
to \emph{prevent information leakage}.  The \textsf{SCX} requirement is the inverse:
we need to \emph{enable auditability} without compromising privacy.

\medskip\noindent
\textbf{Contributions.}
\begin{enumerate}[label=(\arabic*)]
    \item \textbf{Federated Audit $\varepsilon$-Equivalence}
          (Theorem~\ref{thm:fed-equiv}): Under Situs homomorphic composition
          and $\ZK$ verification, $|\Power(T_{\mathrm{fed}}) -
          \Power(T_{\mathrm{full}})| \leq \varepsilon(K, \delta)$,
          with $\varepsilon \to 0$ as $n \to \infty$ and $\delta \to 0$.
          \conditionallyrigorous~/ \openquest~(Situs homomorphism)
    \item \textbf{Information Lower Bound}
          (Theorem~\ref{thm:info-lower}): Without $\ZK$ verification, audit
          power loss is bounded below by a fraction proportional to the
          unverified noise share.  \rigorous
    \item \textbf{Telegraph Public Good Extension}
          (Section~\ref{sec:telegraph}): Decentralized storage commitments
          as a $\ZK$ replacement.  \openquest
\end{enumerate}

% ============================================================
\section{Preliminaries}
% ============================================================

\subsection{SCX Audit Architecture}

For a dataset $\cD = \{(x_i, y_i)\}_{i=1}^n$, the \textbf{Cercis Score}
$S: \cX \to \R_{\geq 0}$ decomposes as
\begin{equation}\label{eq:cercis}
    S(x) = Q(x) + \eta N(x),
\end{equation}
where $Q(x) \in [0,1]$ measures label quality and $N(x) \in [0,1]$ measures novelty.

The \textbf{Situs operator} encodes a distribution into a metric-measure space:
\begin{equation}\label{eq:situs}
    \Situs(P) = (\cX, d_P, \mu_P).
\end{equation}
Two distributions are compared via $\dSitus(P, Q) = \Wass_1(\Situs(P), \Situs(Q))$.

\textbf{Theorem~3} (Noise--Difficulty Unidentifiability). From Cercis Score
observations alone, no algorithm can distinguish label noise from intrinsic
difficulty with power exceeding its significance level.

\subsection{The Federated Audit Setting}

\begin{definition}[Federated Audit Problem]
\label{def:fed-problem}
$K$ parties each hold a private dataset $\cD_k = \{(x_i^{(k)}, y_i^{(k)})\}_{i=1}^{n_k}$.
The auditor $\cA$ can access:
\begin{equation}\label{eq:fed-info}
    \cI_{\mathrm{Fed}} = \{S_k, \nabla S_k, H_{S,k}\}_{k=1}^K \cup \{M_\Pi\},
\end{equation}
where $S_k = \Cercis(\cD_k)$ is the Cercis Score computed locally. The auditor \emph{cannot} access raw $(x, y)$ pairs.
\end{definition}

\subsection{Situs Homomorphic Composition}

\begin{definition}[Situs Homomorphic Composition]
\label{def:situs-hom}
A binary operator $\oplus$ on Situs spaces is \textbf{homomorphic} with respect to dataset union if
\begin{equation}\label{eq:homomorphism}
    \dSitus\!\big(\Situs(\cup_{k=1}^K \cD_k),\;
    \oplus_{k=1}^K \Situs(\cD_k)\big)
    \;\leq\; \phi(n_1, \dots, n_K),
\end{equation}
and $\phi(n_1, \dots, n_K) \to 0$ as $\min_k n_k \to \infty$.
\end{definition}

\subsection{Zero-Knowledge Audit Verification}

\begin{definition}[ZK Audit Proof]
\label{def:zk-proof}
A zero-knowledge proof system for Cercis Score submissions is a triple $(\mathsf{Setup}, \mathcal{P}, \mathcal{V})$ with completeness, soundness, and zero-knowledge properties. The commitment $c_k = \Commit(\cD_k; r_k)$ binds party $k$ to its dataset without revealing it. The $\ZK$ proof $\pi_k$ attests that $S_k$ was correctly computed from $\cD_k$.
\end{definition}

% ============================================================
\section{Main Results}
% ============================================================

\subsection{Federated Audit $\varepsilon$-Equivalence}

\begin{theorem}[Federated Audit Completeness]
\label{thm:fed-equiv}
Under Situs homomorphic composition (Definition~\ref{def:situs-hom}) with error $\phi$, $\ZK$ verification (Definition~\ref{def:zk-proof}), Lipschitz audit regularity, and density regularity, the federated and full-information audit powers satisfy
\begin{equation}\label{eq:fed-equiv-bound}
    \big|\Power(T_{\mathrm{fed}}) - \Power(T_{\mathrm{full}})\big|
    \;\leq\; \varepsilon(K, \delta, \phi),
\end{equation}
where
\begin{equation}\label{eq:epsilon-def}
    \varepsilon(K, \delta, \phi) =
    C_1 \cdot \phi
    + C_2 \cdot \sqrt{\frac{K \cdot \delta}{n_{\mathrm{total}}}}
    + C_3 \cdot \frac{1 + \eta}{\sqrt{\min_k n_k}}
    + C_4 \cdot \negl(\lambda),
\end{equation}
with constants $C_1 = 2LM$, $C_2 = \sqrt{2/\pi}$, $C_3 = \sqrt{2\ln(2K/\xi)}$, $C_4 = 1$.
\end{theorem}

\begin{proof}[Proof sketch]
The proof employs an intermediate-statistic chain decomposition, separately bounding four error sources through triangle inequality:
\begin{enumerate}
    \item \textbf{Situs composition error} ($C_1\cdot\phi$): from Lipschitz continuity of the audit functional.
    \item \textbf{DP distortion error} ($C_2\cdot\sqrt{K\delta/n_{\text{total}}}$): from Gaussian DP trade-off functions.
    \item \textbf{Finite-sample estimation error} ($C_3\cdot(1+\eta)/\sqrt{\min_k n_k}$): from Hoeffding concentration and union bound.
    \item \textbf{Cryptographic soundness error} ($C_4\cdot\negl(\lambda)$): from the ZK proof system's soundness guarantee.
\end{enumerate}
\end{proof}

\subsection{Information Lower Bound Without ZK}

\begin{theorem}[Information Loss Lower Bound]
\label{thm:info-lower}
Without $\ZK$ verification, the audit power loss satisfies
\[
\Power(T_{\mathrm{full}}) - \Power(T_{\mathrm{fed}}) \geq \mathrm{const} \cdot \sum_{k=1}^K \frac{\eta_k |\cD_k|}{|\cD_{\mathrm{total}}|}.
\]
\end{theorem}

This follows from Theorem~3 and the data processing inequality: unverified Cercis Score submissions lose information proportional to the unverified noise share.

\subsection{Telegraph Public Good Extension}
\label{sec:telegraph}

Decentralized storage-layer commitments (e.g., Filecoin, Arweave) can replace $\ZK$ proofs with Proofs of Storage, further tightening the $\varepsilon$ bound. This extension remains an open problem.

% ============================================================
\section{Conclusion}
% ============================================================

FA-Theorem establishes that federated \textsf{SCX} auditing is possible without raw data access, provided the information architecture supplies verifiable audit fingerprints through ZK proofs and Situs homomorphic composition. The $\varepsilon$-equivalence bound characterizes the trade-off among privacy, cryptographic soundness, and statistical accuracy.

% ============================================================
\bibliographystyle{plainnat}
\begin{thebibliography}{10}

\bibitem{scx2024cercis}
SCX Working Group.
\newblock ``The \textsf{SCX} Axiomatic Framework: Cercis, Situs, and Tiresias Operators,''
\newblock \emph{Technical Report}, 2024.

\bibitem{bonawitz2017practical}
K.~Bonawitz et al.
\newblock ``Practical secure aggregation for privacy-preserving machine learning,''
\newblock in \emph{CCS}, 2017.

\bibitem{abadi2016deep}
M.~Abadi et al.
\newblock ``Deep learning with differential privacy,''
\newblock in \emph{CCS}, 2016.

\end{thebibliography}

\end{document}
